\documentclass[11pt,letterpaper]{article}
\usepackage[margin=1in]{geometry}
\usepackage{parskip}
\pagestyle{empty}

\begin{document}

\noindent Krishna Harish, Elkins High School, Missouri City, Texas \hfill \today\\
\noindent krishnaharish2009@gmail.com

\vspace{1em}
\noindent The Editors, \textit{Journal of Chemical Information and Modeling}

\vspace{1em}
\noindent Dear Editors,

Please consider the enclosed manuscript, ``Local Structure--Activity Landscape Roughness Predicts QSAR Error and Keeps Conformal Prediction Intervals Valid on Activity Cliffs,'' for publication as an Article in \textit{JCIM}.

QSAR and deep-learning models guide compound prioritization throughout drug discovery, yet they fail on activity cliffs (pairs of near-identical molecules whose potencies differ by orders of magnitude) without warning. The per-compound reliability tools in everyday use do not catch this failure: ensemble uncertainty ranks ordinary error well but is no better than chance on cliffs, and the applicability domain reports high confidence in exactly the dense regions where cliffs occur. The question a modeler most needs answered, which specific predictions to distrust, is left open for the failure mode that matters most.

This manuscript shows that the missing signal is the local roughness of the structure--activity landscape, computable before assay from a compound's training neighbors without its own activity, and that it is the right per-compound scale on which to judge a prediction's reliability. Across thirty MoleculeACE targets, an activity-free roughness descriptor predicts per-compound error, stays distinct from the applicability domain under partial-correlation and mixed-effects control, and flags activity cliffs before any assay. The practical outcome is a calibrated, cliff-aware conformal predictor: conditioning prediction intervals on roughness keeps 90\% coverage valid across the roughness range and on cliffs, where conditioning on tree variance or the applicability domain leaves it broken. Roughness complements rather than replaces existing uncertainty signals, and a combined score inherits both.

We believe the work fits JCIM and will interest its readership. It reframes activity-cliff reliability as the problem of choosing the right per-compound variable and identifies that variable; it provides a deployable, model-agnostic tool that needs no measured activity and improves on the signals now in use, on precisely the predictions practitioners are most often wrong about; and the findings are reproducible and general, holding across three model classes, neighborhood sizes, and distance metrics, and transferring to aqueous solubility and lipophilicity. All code and the scripts that regenerate every figure and table are openly available, with a citable Zenodo-archived snapshot.

The manuscript is original, is not published or under consideration elsewhere, and as sole author I declare no competing financial interest. Thank you for considering this submission.

\vspace{1em}
\noindent Sincerely,\\[1.5em]
Krishna Harish

\end{document}
